% 原书第177—184页：9.4。
\section{条件独立性与图表示}

因果学习试图从观测数据中重建因果模型。前面已经介绍若干可识别性结果，它们使我们能够从观测分布 $P_{\mathbf X}$ 识别变量集合 $\mathbf X$ 上的结构因果模型的图结构。现在，考虑变量 $\mathbf X=(\mathbf O,\mathbf H)$ 上的结构因果模型 $\mathfrak C$，其中 $\mathbf O$ 是观测变量，$\mathbf H$ 是潜变量。我们仍然可以问：仅凭观测变量的分布 $P_{\mathbf O}$，是否能够识别 $\mathfrak C$ 的图结构？如果能够，应当如何识别？

第7.2.1节讨论了无潜变量时，如何在马尔可夫条件和忠实性假设下学习因果结构的一部分。这两个假设保证d-分离与条件独立性之间存在一一对应关系，因此，可以检验 $P_{\mathbf X}$ 中的条件独立关系，进而重建背后图结构的某些性质。回忆一下，基于独立性的方法原则上在DAG空间中搜索，并输出一个图，或一个图的等价类，使其恰好表示数据中发现的条件独立关系集合。

当因果学习涉及潜变量时，原则上希望在含潜变量的DAG空间中搜索，但这会带来额外困难。我们不知道 $\mathbf H$ 的大小；如果不限制潜变量个数，就必须搜索无穷多个候选图。统计上，这种做法也存在问题：相对于一个DAG满足马尔可夫性和忠实性的分布集合构成弯曲指数族，这为使用BIC等准则提供了依据［Haughton, 1988］；但对于含潜变量的DAG，相应的分布集合并非如此［Geiger and Meek, 1998］。如果直接搜索含潜变量的DAG不可行，能否把每个这样的DAG表示为一个只含观测变量的边缘化图，必要时使用多种边，再搜索这些结构？第9.1节已指出，这种做法也有困难：边缘化图应当取决于原始结构因果模型，仅有原图中包含的信息还不够。如前所述，Bongers et al.［2016］更详细地研究了结构因果模型的边缘化。

因此，本节余下部分稍微转变问题的重点：不再检验一个完整分布是否可能由某种含潜变量的DAG结构诱导，而只考虑特定类型的约束。例如，考虑所有在观测变量 $\mathbf O$ 上满足同一组条件独立关系的分布，这隐含假定了马尔可夫条件与忠实性，然后再问如何表示这组条件独立关系。

一种直接的方案，是假设诱导的分布 $P_{\mathbf O}$ 相对于某个不含潜变量的DAG满足马尔可夫性和忠实性，再像以前一样，输出一类能够表示观测变量分布中条件独立关系的DAG。不过，用DAG表示 $P_{\mathbf O}$ 的条件独立结构有两个众所周知的缺点：（1）在观测变量上用DAG表示条件独立关系集合，可能引起因果误读；（2）独立性模式与DAG中d-分离关系相对应的分布集合，对边缘化并不封闭［Richardson and Spirtes, 2002］。

对于（1），考虑一个诱导分布 $P_{A,B,C,H}$ 的结构因果模型，该分布相对于图9.3左侧的对应DAG满足马尔可夫性和忠实性。在观测分布 $P_{A,B,C}$ 中，唯一能发现的条件或无条件独立关系是 $A\indep C$。因此，图9.3从左起第二个DAG恰好表示这一关系；从这个意义上说，它可以看作PC算法的输出。但其因果解释是错误的：在原始结构因果模型中，对 $C$ 的干预不会影响 $B$，而PC输出的图却暗示 $C$ 对 $B$ 有因果效应。

\begin{figure}[htbp]
\centering
\begin{tikzpicture}[causal,x=1cm,y=1.1cm]
 \node[vtxs] (ch09f3aa) at (0,0) {$A$};
 \node[vtxs] (ch09f3ab) at (1.05,0) {$B$};
 \node[vtxs] (ch09f3ac) at (2.1,0) {$C$};
 \node[latent,minimum size=5.5mm,font=\footnotesize] (ch09f3ah) at (1.55,.95) {$H$};
 \draw[ed] (ch09f3aa)--(ch09f3ab);
 \draw[ed] (ch09f3ah)--(ch09f3ab);
 \draw[ed] (ch09f3ah)--(ch09f3ac);
 \node[align=center,font=\small] at (1.05,-.7) {真实DAG};
 \node[vtxs] (ch09f3ba) at (3.05,0) {$A$};
 \node[vtxs] (ch09f3bb) at (4.1,0) {$B$};
 \node[vtxs] (ch09f3bc) at (5.15,0) {$C$};
 \draw[ed] (ch09f3ba)--(ch09f3bb);
 \draw[ed] (ch09f3bc)--(ch09f3bb);
 \node[align=center,font=\small] at (4.1,-.7) {DAG\\（PC输出）};
 \node[vtxs] (ch09f3ca) at (6.1,0) {$A$};
 \node[vtxs] (ch09f3cb) at (7.15,0) {$B$};
 \node[vtxs] (ch09f3cc) at (8.2,0) {$C$};
 \draw[ed] (ch09f3ca)--(ch09f3cb);
 \draw[bid] (ch09f3cb)--(ch09f3cc);
 \node[align=center,font=\small] at (7.15,-.7) {MAG};
 \node[vtxs] (ch09f3da) at (9.15,0) {$A$};
 \node[vtxs] (ch09f3db) at (10.2,0) {$B$};
 \node[vtxs] (ch09f3dc) at (11.25,0) {$C$};
 \draw[{Circle[open,length=1.5mm]}-{Stealth[length=2.1mm,width=1.5mm]}] (ch09f3da)--(ch09f3db);
 \draw[{Circle[open,length=1.5mm]}-{Stealth[length=2.1mm,width=1.5mm]}] (ch09f3dc)--(ch09f3db);
 \node[align=center,font=\small] at (10.2,-.7) {PAG\\（修改后PC或FCI输出）};
\end{tikzpicture}
\caption{从左侧的结构因果模型出发，右侧三个图都编码了条件独立关系集合 $A\indep C$。由于会导致错误的因果解释，DAG不宜作为这里因果学习方法的输出。在这个例子中，IPG与潜在投影所得的ADMG均与MAG相同。}
\label{fig:ch09-3}
\end{figure}

\begin{figure}[htbp]
\centering
\begin{tikzpicture}[causal,x=1.6cm,y=1.3cm]
 \node[vtx] (ch09f4a) at (0,0) {$A$};
 \node[vtx] (ch09f4b) at (1.5,0) {$B$};
 \node[vtx] (ch09f4c) at (3,0) {$C$};
 \node[vtx] (ch09f4d) at (4.5,0) {$D$};
 \node[latent] (ch09f4h) at (2.5,1.2) {$H$};
 \draw[ed] (ch09f4a)--(ch09f4b);
 \draw[ed] (ch09f4c)--(ch09f4d);
 \draw[ed] (ch09f4h)--(ch09f4b);
 \draw[ed] (ch09f4h)--(ch09f4d);
\end{tikzpicture}
\caption{例子取自 Richardson and Spirtes［2002，图2（i）］，说明DAG对边缘化并不封闭。不存在一个节点集为 $\mathbf O=\{A,B,C,D\}$ 的DAG，能够编码含 $H$ 的原图中的全部条件独立关系。}
\label{fig:ch09-4}
\end{figure}


对于（2），图9.4的例子取自 Richardson and Spirtes［2002］。它显示了变量 $\mathbf X=(\mathbf O,\mathbf H)$ 上一个结构因果模型的结构，其分布相对于DAG $\mathcal G$ 满足马尔可夫性和忠实性，即 $\mathcal G$ 表示 $\mathbf X$ 中的全部条件独立关系。这个模型有如下性质：不存在任何定义在 $\mathbf O$ 上的DAG，能够表示 $P_{\mathbf O}$ 中的条件独立关系。正是在这个意义上，DAG对边缘化并不封闭。

下一小节讨论用观测变量 $\mathbf O$ 上的图表示条件独立性的一些方案。但这些图未必具有直观的因果含义，可能很难从它们推断出变量 $\mathbf X=(\mathbf O,\mathbf H)$ 上原始结构因果模型的结构性质。例如，第6.6节那样用于协变量调整的图判据，必须针对每一种图重新提出并证明。

\subsection{图}

前面，我们用图表示结构因果模型的结构关系，见定义3.1和6.2。本节的目的有所不同：用图表示结构因果模型所诱导分布中的约束。第9.4节主要考虑条件独立关系；其他约束将在第9.5节详细讨论。我们已经看到，存在潜变量时，用DAG表示条件独立关系并不理想。这些不足促使因果推断研究者发展新的图表示。例如，Richardson and Spirtes［2002］引入了\textbf{极大祖先图（maximal ancestral graphs，MAG）}，并证明它们构成对边缘化封闭的最小DAG超类，见前面的讨论。MAG是含有有向边和双向边的混合图。\footnote{事实上，MAG还可以含有无向边，因此也能对选择偏倚建模。详细讨论见 Richardson and Spirtes［2002］。}
它采用稍有不同的分离判据：以\textbf{m-分离（m-separation）}代替d-分离［Richardson and Spirtes, 2002］。于是，对于每一个含潜变量的DAG，在观测变量上都有唯一的MAG，通过m-分离表示同一组条件独立关系。Richardson and Spirtes［2002，第4.2.1节］给出了简单的构造程序，图9.3展示了一个例子。

这种映射不是一一对应的。每个MAG都能由无穷多个不同的DAG构造出来，这些DAG可以包含任意多个潜变量。与DAG类似，马尔可夫条件把MAG中的图分离关系与条件独立性联系起来。表示同一组m-分离关系的不同MAG，归入同一个马尔可夫等价类［Zhang, 2008b］；这个等价类本身通常用\textbf{部分祖先图（partially ancestral graph，PAG）}表示，概览见表9.1。在PAG中，边的端点可以是圆圈，表示该端既可能为箭头，也可能为箭尾，见图9.3。Ali et al.［2009］给出了判断两个MAG是否马尔可夫等价的图判据。

% 原书表9.1；由第9.4.1节输入，非第9.5节正文。
\begin{table}[htbp]
\centering
\small
\setlength{\tabcolsep}{4pt}
\begin{tabularx}{\linewidth}{p{4.0cm}*{4}{>{\centering\arraybackslash}X}}
\toprule
图对象&DAG（无潜变量）&MAG&IPG&ADMG（含嵌套马尔可夫模型）\\
\midrule
边类型：有向/无向/双向/组合
 &$\checkmark/-/-/-$&$\checkmark/\checkmark/\checkmark/-$&$\checkmark/-/\checkmark/-$&$\checkmark/-/\checkmark/\checkmark$\\
正确的因果解释&$\times$&$\checkmark$&$\checkmark$&$\checkmark$\\
全局马尔可夫性的图分离&d-分离&m-分离&m-分离&m-分离\\
有效调整集的判据&$\checkmark$&$\checkmark$&?&$\checkmark$\\
识别干预分布的算法&$\checkmark$&?&?&$\checkmark$\\
等价类的表示&CPDAG（马尔可夫）&PAG（马尔可夫）&POIPG（马尔可夫）&?（嵌套马尔可夫）\\
基于独立性的学习方法&PC, IC, SGS&FCI&FCI&--\\
基于得分的学习方法&GDS, GES&适用于线性/二值/离散SCM&?&适用于二值/离散SCM\\
能否编码全部等式约束&$\times$&$\times$&$\times$&$\checkmark$（观测变量离散时）\\
能否编码全部约束&$\times$&$\times$&$\times$&$\times$\\
\bottomrule
\end{tabularx}
\caption{考虑一个结构因果模型，其（已观测）变量为$\mathbf O$、（潜）变量为$\mathbf H$，并诱导出分布$P_{\mathbf O,\mathbf V}$。应如何对观测分布$P_{\mathbf O}$建模？我们希望采用允许任意多个潜变量的结构因果模型，但这一模型类的性质不利于因果学习。本表概括一些可替代的模型类；当前研究尤其关注MAG与ADMG。}
\end{table}


\begin{example}[Y结构]
即使单个MAG，也能表示任意多个潜变量，因此，人们可能会怀疑：由含潜变量的DAG构造的PAG，是否还能包含任何非平凡的因果信息？例如，图9.3的PAG没有说明，究竟存在一条从 $C$ 到 $B$ 的有向路径，还是存在一个潜变量，分别沿有向路径指向 $C,B$。图9.2右图展示了Y结构，其中 $Z_1,Z_2,Y$ 两两之间没有直接连接。现在，考虑任意多个变量上的结构因果模型，它包含 $X,Z_1,Z_2,Y$ 四个变量，并在这四个变量上诱导出与Y结构相同的条件独立关系。那么，可以断定：对应PAG中含有有向边 $X\to Y$。此外，$X,Y$ 之间的因果关系必定没有混杂［例如 Mani et al., 2006；Spirtes et al., 2000，图7.23］。在任何结构因果模型中，只要 $X,Y$ 存在混杂，或者 $X$ 不是 $Y$ 的祖先，就会产生不同的条件独立关系集合。
\end{example}

前面已提到，MAG等图对象主要是为表示条件独立关系而构造的，并不是为了直观展示结构因果模型；后者才是定义3.1引入图时的用途。因此，它们的因果语义更加复杂。例如，在MAG中，边 $A\to B$ 表示：在包含潜变量的原始DAG中，$A$ 是 $B$ 的祖先，而 $B$ 不是 $A$ 的祖先，即祖先关系得以保留。图9.3的PAG应解释为：“在原始DAG中，可能存在从 $C$ 到 $B$ 的有向路径，也可能存在潜在共同原因，还可能二者兼有。”因此，在这类图中进行因果推理，即计算干预分布，也更加复杂［例如 Spirtes et al., 2000；Zhang, 2008b］。Perkovic et al.［2015］刻画了有效调整集，见第6.6节，其结论既适用于DAG，也适用于MAG。

除了MAG与PAG，还可以使用\textbf{诱导路径图（induced path graphs，IPG）}，以及表示一组IPG的\textbf{完备部分定向诱导路径图（completed partially oriented induced path graphs，POIPG）}［Spirtes et al., 2000，第6.6节］。这类图最初用于表示快速因果推断算法的输出，见第9.4.2节。考虑一个相对于某MAG满足马尔可夫性与忠实性的分布。每个MAG都是IPG，反之则未必，因此，该MAG的马尔可夫等价类包含在相应IPG的马尔可夫等价类中。由此，PAG通常比POIPG包含更多因果信息［Zhang, 2008b，附录A］。

另一种方法是从含潜变量的原始DAG出发，施加\textbf{潜在投影（latent projection）}［Pearl, 2009，定义2.6.1；Verma and Pearl, 1991，其中称为“嵌入模式”］。这一操作以含观测变量和潜变量的图 $\mathcal G$ 为输入，在观测变量上构造新图 $\widetilde{\mathcal G}$。精确定义可见 Shpitser et al.［2014，定义4］。所得结构称为\textbf{无环有向混合图（acyclic directed mixed graph，ADMG）}，包含有向边和双向边。m-分离同样导出一种马尔可夫性［Richardson, 2003］。于是，可以搜索ADMG空间，代替搜索含潜变量的DAG空间。

第9.5节将说明，含潜变量的DAG在观测变量上的分布，还满足条件独立性以外的约束。ADMG可按下面的方法纳入其中一些约束。基本思想是定义\textbf{嵌套马尔可夫性（nested Markov property）}［Richardson et al., 2012, 2017；Shpitser et al., 2014］：如果一个分布不仅满足图结构蕴含的某些条件独立关系，还满足其他约束，就称其相对于ADMG具有嵌套马尔可夫性，例子见第9.5.1节。即便如此，嵌套马尔可夫性也无法编码全部约束；不过，在离散情形下，它确实能编码全部等式约束［Evans, 2015］。因此，有下面的包含关系［Shpitser et al., 2014］：\chninenote{此处集合的第一行，以及表9.1图注中的原书记号均为 $P_{\mathbf O,\mathbf V}$，但本节一直将潜变量集合记作 $\mathbf H$；其中 $\mathbf V$ 应理解为 $\mathbf H$。}
\[
\begin{aligned}
 &\{P_{\mathbf O}:P_{\mathbf O,\mathbf V}\text{ 由含潜变量的DAG }\mathcal G\text{ 诱导}\}\\
 &\quad\subseteq\{P_{\mathbf O}:P_{\mathbf O}\text{ 相对于对应ADMG具有嵌套马尔可夫性}\}\\
 &\quad\subseteq\{P_{\mathbf O}:P_{\mathbf O}\text{ 相对于对应ADMG具有马尔可夫性}\}.
\end{aligned}
\]


对于离散数据和满足通常马尔可夫性的ADMG，Evans and Richardson［2014］给出了参数化形式。该参数化可以推广到嵌套马尔可夫模型，并用于计算受约束的极大似然估计［Shpitser et al., 2012］。如果任意一对节点之间至多有一种边，就称ADMG为\textbf{无弓图（bow-free）}。在线性高斯模型中，这一子类允许参数可识别［Brito and Pearl, 2002a］；此外，已有算法可计算极大似然估计［Drton et al., 2009a］，或用于因果学习［Nowzohour et al., 2015］。

\textbf{链图（chain graphs）}由有向边和无向边构成，且不允许部分有向环［Lauritzen, 1996，第2.1.1节］。关于链图已有大量研究；综述见 Lauritzen［1996］，因果解释见 Lauritzen and Richardson［2002］。注意，对链图已有不同的马尔可夫性定义［Lauritzen and Wermuth, 1989；Frydenberg, 1990；Andersson et al., 2001］。

由此可见，用图表示约束，尤其是在存在潜变量时，并非简单的问题，至今仍是活跃的研究领域；目前为止，主要讨论的是条件独立性。Sadeghi and Lauritzen［2014］联系了若干类混合图，并讨论其马尔可夫性。通常，这些图对象及其分离判据相当复杂，将边与因果效应的存在性联系起来也并不容易，尽管嵌套马尔可夫模型或许可视为朝简化方向迈出的一步。令人惊讶的是，尽管存在这些困难，在某些情况下，例如例9.5的Y结构，我们仍然能够学习因果祖先关系。

\subsection{快速因果推断}

已经看到，对于结构学习而言，PAG可能比CPDAG更适合作为输出。事实上，可以修改PC算法，使其输出PAG［Spirtes et al., 2000，第6.2节］。虽然这种简单修改在许多例子中有效，但一般而言并不正确。PC算法的每一次迭代都会考虑一对当前相邻节点，例如 $A,B$，再搜索使它们d-分离的集合。为了显著提高速度，根据第7.2.1节的引理7.8（ii），搜索只考虑 $A,B$ 当前邻居集合的子集。但是，存在潜变量时，仅搜索邻居集合的子集就不再充分［Verma and Pearl, 1991，引理3］。Spirtes et al.［2000，第6.3节］给出了一个例子，修改后的PC算法在其中无法找到d-分离集。

\textbf{快速因果推断（fast causal inference，FCI）}算法［Spirtes et al., 2000］解决了这一问题，输出表示一组MAG的PAG。Zhang and Spirtes［2005］与 Zhang［2008a］证明，对原始FCI算法作少许修改即可达到完备性，即使输出包含尽可能多的信息。如果条件独立关系来自某个含潜变量的DAG，算法输出就确实表示正确的对应PAG。

若干FCI变体能显著提速。Spirtes［2001］提出限制条件集合的大小，称为anytime FCI；Colombo et al.［2012］同时减少条件独立性检验的次数和条件集合的大小，称为真正快速因果推断（really fast causal inference）。两者包含的信息都可能略少于FCI。后续提出的FCI+则兼具速度和完备性［Claassen et al., 2013］。

另一条途径是对MAG或MAG的等价类评分。目前，这类评分函数仅针对某些结构因果模型类别存在，例如线性结构因果模型［Richardson and Spirtes, 2002］；而且，据我们所知，还没有高效搜索MAG空间的方法［Mani et al., 2006］。Silva and Ghahramani［2009］讨论了一种学习混合图的贝叶斯方法。
\FloatBarrier
